% paper/sections/10e_verification_summary.tex
%
% §10.4  検証結果のまとめ
% ═══════════════════════════════════════════════════════════════════════════════
\subsection{検証結果のまとめ}
\label{sec:verify_summary}
% ═══════════════════════════════════════════════════════════════════════════════

表~\ref{tab:verification_summary} に各コンポーネントの検証結果を総括する．
主要な知見を以下に整理する．

\begin{enumerate}
  \item \textbf{空間離散化：}
        CCD 微分は周期境界テストで $\Ord{h^{6.0}}$ を達成し，設計精度を確認した．
        DCCD フィルタは Nyquist モードを完全に抑制する（$H(\pi)=0$）．
  \item \textbf{曲率計算：}
        $\phi$ 直接微分（経路 A）では $\Ord{h^5}$ の高次収束を示すが，
        本番経路 B（$\psi \to \phi$ ロジット逆変換）は界面厚み $\varepsilon = \Ord{h}$ に律速される．
        ただし GFM 界面処理精度（$\Ord{h^2}$）が全体を律速するため，
        曲率計算経路の選択は実用上の精度に影響しない．
  \item \textbf{PPE 求解：}
        欠陥補正法は $k=3$ 回の反復で CCD 演算子の $\Ord{h^6}$ 精度を完全に発現させ，
        LGMRES と同等の精度を達成する．
        FD 直接法と比較して $N=128$ で $4{,}900$ 万倍高精度であり，コスト増は $3.4$ 倍に過ぎない．
  \item \textbf{時間積分：}
        TVD-RK3 は ODE テストで $\Ord{\Delta t^{3.00}}$ の精確な 3 次収束を確認した．
\end{enumerate}

\begin{table}[htbp]
\centering
\caption{計算コンポーネント検証結果の総括}
\label{tab:verification_summary}
\begin{tabular}{llll}
\toprule
コンポーネント & 期待精度 & 実測精度 & 備考 \\
\midrule
CCD 微分（周期） & $\Ord{h^6}$ & $\Ord{h^{6.0}}$ & $N = 256$ で丸め誤差到達 \\
DCCD フィルタ   & $H(\pi) = 0$ & $H(\pi) = 0$ & チェッカーボード完全抑制 \\
曲率（正弦波，$\phi$ 直接） & $\Ord{h^6}$ & $\Ord{h^{5.0}}$ & CCD 微分精度のベースライン \\
CLS 質量保存    & 保存的     & $\Ord{10^{-5}}$ ($N=256$) & 形状誤差は $\varepsilon$ 律速 \\
GFM Laplace 圧力跳び & 安定収束 & $\Ord{h^1}$ & 粗格子でも安定；$N=128$ で誤差 $1.8\%$ \\
PPE（欠陥補正 $k{=}3$） & $\Ord{h^6}$ & $\Ord{h^{7.0}}$ & LGMRES と同等；ディリクレ超収束 \\
TVD-RK3         & $\Ord{\Delta t^3}$ & $\Ord{\Delta t^{3.00}}$ & ODE テストで確認 \\
\bottomrule
\end{tabular}
\end{table}
